\documentclass[10pt]{article}
\usepackage[T1]{fontenc}
\usepackage[utf8]{inputenc}
\usepackage[french]{babel}
\usepackage{amsmath,amssymb,mathrsfs,mathpazo}
\usepackage{geometry}
\geometry{a4paper,margin=2cm}
\usepackage{graphicx,array,tabularx,tabularray,colortbl,multirow,multicol,booktabs,enumitem,pifont,ulem,fancybox,diagbox,cancel}
\usepackage{pstricks,pstricks-add,pst-plot,pst-tree,pst-node,pst-text,pst-func,pst-math,pst-bspline,pst-3dplot,pst-eucl,pst-solides3d}
\let\pscalculate\undefined
\usepackage{pst-bezier}
\usepackage{tikz,pgfplots,tkz-tab}
\usetikzlibrary{arrows,arrows.meta,patterns,calc,positioning,trees,decorations.markings,shapes.geometric,shapes.misc,shapes.arrows,intersections,decorations.pathmorphing,decorations.pathreplacing,backgrounds,shadows,babel}
\pgfplotsset{compat=1.18}
\usepgfplotslibrary{fillbetween}
\usepackage[np]{numprint}
\usepackage{listings,makecell,pas-tableur,twoopt}
\let\barre\undefined
\usepackage[pstricks]{tablvar}
\definecolor{darkgreen}{RGB}{0,100,0}
\providecommand{\red}{\color{red}}
\providecommand{\blue}{\color{blue}}
\providecommand{\green}{\color{green}}
\providecommand{\black}{\color{black}}
\providecommand{\white}{\color{white}}
\providecommand{\gray}{\color{gray}}
\providecommand{\cyan}{\color{cyan}}
\providecommand{\magenta}{\color{magenta}}
\providecommand{\yellow}{\color{yellow}}
\providecommand{\violet}{\color{violet}}
\providecommand{\orange}{\color{orange}}
\providecommand{\euro}{\texteuro}
\providecommand{\EUR}{\texteuro}
\providecommand{\ieme}{\textsuperscript{e}}
\providecommand{\degres}{\ensuremath{^\circ}}
\providecommand{\degree}{\ensuremath{^\circ}}
\providecommand{\vv}[1]{\overrightarrow{#1}}
\providecommand{\Coord}[1]{#1}
\providecommand{\decosix}{}
\providecommand{\starredbullet}{\ensuremath{\bullet}}
\pagestyle{empty}
\setlength{\parindent}{0pt}
\usepackage{adjustbox}\DeclareUnicodeCharacter{2236}{:}\DeclareUnicodeCharacter{2190}{\ensuremath{\leftarrow}}
\providecommand{\R}{\mathbb{R}}
\providecommand{\N}{\mathbb{N}}
\providecommand{\Z}{\mathbb{Z}}
\providecommand{\Q}{\mathbb{Q}}
\providecommand{\D}{\mathbb{D}}
\providecommand{\C}{\mathbb{C}}
\providecommand{\e}{{\mathrm{e}}}
\providecommand{\vect}[1]{\overrightarrow{#1}}
\providecommand{\vectt}[1]{\overrightarrow{\mathrm{#1}}}
\providecommand{\vvt}[1]{\overrightarrow{\mathrm{#1}}}
\providecommand{\np}[2][]{{#2~#1}}
\providecommand{\vv}[1]{{\overrightarrow{#1}}}
\renewcommand{\d}{{\mathrm{d}}}
\providecommand{\degres}{{^\circ}}
\providecommand{\degre}{{^\circ}}
\providecommand{\blue}{\color{blue}}
\providecommand{\red}{\color{red}}
\providecommand{\enskip}{\quad}
\providecommand{\llbracket}{\lbrack\!\lbrack}
\providecommand{\rrbracket}{\rbrack\!\rbrack}
\providecommand{\euro}{{\text{€}}}
\providecommand{\no}{{\text{n°}}}
\providecommand{\interval}[2]{{[#1;#2]}}
\providecommand{\intervFO}[2]{{[#1;#2[}}
\providecommand{\dotfill}{\ldots}

\providecommand{\ds}{\displaystyle}
\providecommand{\textding}[1]{#1}
\providecommand{\Oij}{\(\mathrm{O};\vec{\imath},\vec{\jmath}\)}
\providecommand{\Oijk}{\(\mathrm{O};\vec{\imath},\vec{\jmath},\vec{k}\)}
\providecommand{\Ouv}{\(\mathrm{O};\vec{u},\vec{v}\)}
\providecommand{\intervFF}[2]{\ensuremath{[#1;#2]}}
\renewcommand{\euro}{\ifmmode\text{\texteuro}\else\texteuro\fi}
\DeclareUnicodeCharacter{00B0}{\ifmmode^{\circ}\else\textdegree\fi}
\DeclareUnicodeCharacter{20AC}{\ifmmode\text{\texteuro}\else\texteuro\fi}
% Énoncés : Éducation nationale. Transcription et corrigés : APMEP, Denis Vergès et contributeurs. Découpage et classement : CoopMaths, Sébastien Lozano, Éric Elter et contributeurs. Adaptation de la présentation : Maths Explorer.
% CC BY-SA : https://coopmaths.fr/www/about/
% https://coopmaths.fr/alea/static/bac/2023/tex/bac_2023_09_sujet1_metropole_3_cor.tex
\begin{document}
\noindent\textit{APMEP / CoopMaths. CC BY-SA. Présentation adaptée par Maths Explorer.}\par\medskip
\medskip


\emph{Les probabilités demandées seront données à $10^{-3}$ près.}






















\medskip

\textbf{Partie A}

\medskip

\begin{enumerate}
\item 
D'après l'énoncé :

\begin{itemize}
\item[$\bullet~~$] $p_A(T) = 0,97$ ;
\item[$\bullet~~$] $p_{\overline{A}}\left(\overline{T}\right) = 0,957$ ;
\item[$\bullet~~$] $p(T) = 0,2$.
\end{itemize}

D'où l'arbre pondéré :

\begin{center}
\pstree[treemode=R,nodesepA=0pt,nodesepB=2.5pt,treesep = 1cm,levelsep=2.5cm]{\TR{}}
{\pstree{\TR{$A$~}\taput{$x$}}
	{\TR{$T$}\taput{0,97}
	\TR{$\overline{T}$}\tbput{0,03}
}
\pstree{\TR{$\overline{A}$~}\tbput{$1 - x$}}
{\TR{$T$}\taput{0,043}
\TR{$\overline{T}$}\tbput{0,957}
}
}
\end{center}

\item 
	\begin{enumerate}
		\item 
D'après la loi des probabilités totales :
		
$p(T) = p(A \cap T) + p\left(\overline{A} \cap T\right) = p(A) \times p_{A}(T) + p\left(\overline{A}\right) \times p_{\overline{A}}(T) = x \times 0,97 + (1 - x) \times 0,043 = 0,97x + 0,043 -0,043x = 0,927x + 0,043$.
		\item 
Comme $p(T) = 0,2 = 0,927x + 0,043 \iff 0,157 = 0,927x \iff \dfrac{0,157}{0,927} = x$.
		
Or $\dfrac{0,157}{0,927} \approx \np{0,1694}$ soit 0,169 au millième près.

$x  = p(A) \approx 0,169$.
	\end{enumerate}	
\item 


L'affirmation se traduit par : $p_T(A) > 0,8$.

Or $p_T(A) = \dfrac{p(T \cap A)}{p(T)} = \dfrac{p(A \cap T)}{p(T)} = \dfrac{p(A) \times p_A(T)}{p(T)} \approx \dfrac{0,169 \times 0,97}{0,2}$, soit 

$p_T(A) \approx \dfrac{\np{0,16393}}{0,2} \approx \np{0,81965}$, soit environ 81,97\,\% : l'affirmation est vraie.
\end{enumerate}





\textbf{Partie B}

\medskip









\begin{enumerate}
\item 


Les tirages successifs étant indépendants et chaque personne ayant un probabilité d'être allergique égale à 0,08, $X$ suit une loi de Bernoulli de paramètres $n = 150$ et $p = 0,08$.
\item 
La calculatrice donne $p(X = 20) \approx \np{0,00820}$, soit 0,008 au millième près.
\item 
La calculatrice donne $p(X \leqslant 14) \approx \np{0,7797}$, donc $p(X \geqslant 15) = 1 - p(X \leqslant 14) \approx 1 - \np{0,7797} \approx \np{0,2203}$ soit 0,220 au millième près.
\end{enumerate}

\bigskip

\end{document}
